\section{Related work}
\label{sec:related}

Self-correction studies ask whether models can fix errors when told to look for them, and largely answer no. Models asked to review their own answers rarely improve them \citep{huang2024}, and they struggle to locate errors even when told one exists \citep{tyen2024}. They fail to improve plans by self-critique \citep{valmeekam2023selfcritique,stechly2024selfverification}, and they remain blind to injected errors in their own generations even in benchmarks built to elicit correction \citep{tsui2025}. Prompted-critique studies find the same leniency toward supplied reasoning: models grading step-by-step solutions accept a large majority of wrong steps \citep{reflect2026,vair2026}, including when a single step was surgically edited \citep{breakingchain2026}. A related line perturbs reasoning to measure dependence rather than detection. Corrupting a chain of thought to see whether the final answer changes tests whether the model uses its stated reasoning \citep{lanham2023}, and invalid reasoning in few-shot demonstrations barely hurts downstream accuracy \citep{wang2022,schaeffer2023invalid}. Faithfulness studies likewise find that stated reasoning can diverge from the computation producing the answer \citep{turpin2023unfaithful,arcuschin2025wild,chen2025unfaithful}, and prompted verification pipelines reduce errors when the model is told to check \citep{madaan2023selfrefine,dhuliawala2024cove,ling2023deductive}. Our question is upstream of all of these. It is not whether a model can find an error when asked, nor whether its answer depends on its stated reasoning. It is whether the model notices an error it was never told about, in reasoning it believes is its own.

Error propagation gives the question its stakes. Hallucinations snowball into confident elaboration \citep{zhang2023snowball}, and single wrong steps compound in compositional tasks \citep{dziri2023} and in next-token training itself \citep{bachmann2024pitfalls}. Reasoning accuracy also moves under surface perturbations of the problem: irrelevant context \citep{shi2023distracted}, premise order \citep{chen2024premise}, and entity or number substitutions \citep{mirzadeh2025gsmsymbolic}. That line perturbs the problem and watches the answer; we perturb the model's own solution and watch whether the model notices.

The nearest work to ours measures unprompted checking, and we differ from it on the variable that turns out to matter. The Self-Verification Dilemma \citep{svd2026} finds that models rarely re-verify steps in naturally produced traces, framed as a general token-budget tradeoff, on open models. We plant audited errors rather than relying on natural ones, which fixes the ground truth. We separate readable from computed content, which locates the failure precisely. And we test frontier systems, which shows the readable side improving while the computed side does not. A cross-conversation analog \citep{reclaim2026} finds models re-derive conclusions rather than trusting remembered ones across sessions; our result is the within-trace complement, where the same deference that is prudent across sessions becomes total within one.

Concurrent work reaches the deference conclusion by other routes. Self-Correction Bench attributes uncorrected errors to a capability that exists but is not activated, and the Self-Correction Illusion shows that role and addressability manipulations gate correction \citep{tsui2025,sci2026}. Our contribution to this converging picture is localization: the dissociation is specific to committed computed values, survives provenance and instruction manipulations, and persists when the capability is demonstrably deployable in the same context.

Yee et al.\ \citep{yee2024faithful} report that models frequently recover from corrupted intermediate values in GSM8K reasoning chains, in apparent contradiction with the wall we document. Two factors close the gap. First, part of the reported recovery reflects stimulus memorization rather than verification. Re-running their protocol with renumbered problems, in the style of GSM-Symbolic \citep{mirzadeh2025gsmsymbolic}, GPT-4o's recovery falls from 66\% to 52\% (paired $p=0.004$), with roughly a quarter of its original recoveries flipping to error propagation. On those items the model was restoring remembered numbers, not recomputing (Appendix~\ref{app:renumber}). Second, and more fundamentally, the two designs probe different moments of commitment. Their truncation cuts immediately after the corrupted digits, inside an unfinished calculator annotation, so the model resumes mid-expression---and completing the expression entails recomputing it. Our injections sit after completed statements. Consistent with this distinction, recovery in their own propagated condition, where the corrupted value additionally appears in earlier finished text, falls to roughly a third. We read their results as characterizing uncommitted, in-flight values, and ours as characterizing committed ones.

A mechanistic line of work converges on our deference reading from inside the network. Probing studies find that models internally represent whether an arithmetic step is correct even while their behavior ignores it \citep{validationgap2025}. This matches what we observe behaviorally: the failure is not missing information but an unused check. The deference reading parallels sycophancy, where models yield to a user's stated view against their own knowledge \citep{sharma2024sycophancy}; here the authority deferred to is computed-looking text. Relatedly, models executing long programs degrade as traces lengthen \citep{execdecay2026}; we hold length fixed and vary only the computation a check requires, isolating verification cost from execution burden. The readable side of our contrast is a retrieval ability with known context-position effects \citep{liu2024lost}; our planted contradictions sit adjacent to the site, the easy case.

Trained verifiers are the deliberate foil for our result. Process reward models supervise individual reasoning steps and readily learn to catch computation errors \citep{cobbe2021verifiers,uesato2022process,lightman2023,wang2024mathshepherd}. Trained critics catch bugs that reviewers miss \citep{saunders2022selfcritique,mcaleese2024critics}, and models can be trained to revise their own output \citep{welleck2023selfcorrect}. The checking ability is plainly learnable. Our finding is that the generating model does not deploy it unprompted, at any scale we test, and that neither instruction nor demonstration turns it on at inference time. The gap between what a dedicated verifier catches and what the generator checks by default is exactly the surface our measurement exposes.

Finally, our elicitation borrows from work on forced continuation, where models are made to continue from imposed prefixes to study safety recovery \citep{qi2025,vonrecum2026}. We use the same mechanism for measurement rather than attack. Continuation from the model's own perturbed trace is what lets us observe the unprompted default, and error compounding in multi-step reasoning \citep{dziri2023} is why that default matters downstream.
